\section{路由实验：研究质量与预算分配}\label{sec:routing}
\subsection{四种组合的最终结果}
路由短批在 ahc026 上组合两类动作路由与两类模型路由。各组从相同初始程序开始，研究循环上限为 360 秒、总 HTTP 请求上限为 20 次，候选上限为 10 个。模型调用、候选验证和主动结束共同决定实际使用的预算。

\reportfigure{routing_tradeoff.pdf}{四种路由组合的最终质量与单轮总耗时。质量为最终集相对初始程序的得分提升；总耗时包含最终测试。每组对应单次独立运行（$n=1$）。}{fig:tradeoff}

\begin{smalltable}
\begin{tabularx}{\linewidth}{@{}Lrrrr@{}}\toprule
路由组合（动作 / 模型） & 最终得分 & 相对 $c_0$ & 总耗时（秒） & 较基线分差\\\midrule
基线路由（LLM / 规则） & 139736 & 5.970\% & 392.039 & 0\\
动作决策（Decisions / 规则） & 139972 & 6.149\% & 389.380 & +236\\
模型决策（LLM / Decisions） & 139559 & 5.836\% & 389.913 & $-177$\\
双决策（Decisions / Decisions） & 132540 & 0.513\% & 112.910 & $-7196$\\\bottomrule
\end{tabularx}
\end{smalltable}

前三组的总时长接近，最终得分落在 139559--139972 之间。动作决策组的得分最高，比基线路由高 236 分，增幅 0.169\%。双决策组在 112.910 秒结束，较基线路由减少 71.199\% 的单轮总耗时，最终改进幅度为 0.513\%。它体现了一种更早停止、较少使用强生成模型的预算分配方式。

\subsection{把原生分数还原为搬运成本}
ahc026 要求按编号取出箱子，移动一段堆叠箱子时付出与移动数量相关的能耗。题目对单例使用 $S=\max(1,10000-V)$ 计分，其中 $V$ 为箱子搬运能耗。本批最终用例的分数均高于最低计分值，因此 15 例的总能耗可由 $V_{\mathrm{sum}}=150000-S_{\mathrm{sum}}$ 得到。

\begin{smalltable}
\begin{tabularx}{\linewidth}{@{}Lrr@{}}\toprule
方案 & 15 例总搬运能耗 & 较初始程序减少\\\midrule
初始程序 $c_0$ & 18136 & ---\\
基线路由 & 10264 & 43.405\%\\
动作决策 & 10028 & 44.707\%\\
模型决策 & 10441 & 42.429\%\\
双决策 & 17460 & 3.727\%\\\bottomrule
\end{tabularx}
\end{smalltable}

这项换算赋予分数变化更直观的任务意义：动作决策组把累计搬运能耗从 18136 降至 10028，主要减少了为取出下一只箱子而进行的额外搬移。

\subsection{模型升级与停止行为}
基线路由、动作决策与模型决策三组都发生了一次从 Qwen 到 Sol 的升级；双决策组全程使用 Qwen。前两组在接近时间上限时结束；模型决策与双决策由代理主动结束。四组生成候选分别为 5、7、6、6 个，全部通过开发验证，最终各交付一个通过 15 例测试的程序。

动作决策组将较多生成机会留给已经升级的 Sol，完成了五次 Sol 生成请求。双决策的 19 次 HTTP 请求中有 13 次用于动作或模型决策，余下 6 次由 Qwen 执行。统一请求预算下，决策频率会直接影响剩余的生成机会。模型选择和停止策略需要共同考虑已经获得的质量、剩余请求和下一次实验的预期收益。

基线路由组有 1 次 Sol 请求因超时（timeout）失败并计入请求额度。鉴于各组仅运行单次，后续仍需通过重复实验剥离服务延迟波动、采样随机性与路由策略各自的贡献。
\source{[S3] 冻结组合、逐例最终得分和调用记录；搬运能耗公式来自归档任务说明，初始逐例结果来自 [S2] 的 c0 交付轮。}
